\begin{lemma}
\label{lemma-regular-koszul-regular}
Let $R$ be a ring, $M$ an $R$-module, and $f_1, \ldots, f_r \in R$
such that for $i = 1, \ldots, r$ multiplication by $f_i$
is injective on $M/(f_1, \ldots, f_{i - 1})M$. Then $f_1, \ldots, f_r$
is $M$-Koszul regular. In particular, an $M$-regular sequence
is $M$-Koszul-regular and any regular sequence is Koszul-regular.
\end{lemma}

\begin{proof}
Let $R$, $M$, $f_1, \ldots, f_r$ be as in the first sentence of the lemma.
If $r = 1$, it is immediate that $f_1$ is $M$-Koszul-regular.
Assume $r > 1$. Since $f_1$ is a nonzerodivisor on $M$, we obtain
a short exact sequence of complexes:
$$
0 \to K_\bullet(f_2, \ldots, f_r) \otimes M
\xrightarrow{f_1}
K_\bullet(f_2, \ldots, f_r) \otimes M \to
K_\bullet(\overline{f}_2, \ldots, \overline{f}_r) \otimes M/f_1M \to 0
$$
Here $\overline{f}_i$ is the image of $f_i$ in $R/(f_1)$. By
Lemma \ref{lemma-cone-koszul}
the complex $K_\bullet(f_1, \ldots, f_r)$
is isomorphic to the cone of multiplication by $f_1$
on $K_\bullet(f_2, \ldots, f_r)$. Thus
$K_\bullet(R, f_1, \ldots, f_r) \otimes M$ is isomorphic
to the cone on the first map. Hence
$K_\bullet(\overline{f}_2, \ldots, \overline{f}_r) \otimes M/f_1M$
is quasi-isomorphic to $K_\bullet(f_1, \ldots, f_r) \otimes M$.
As $R/(f_1)$, $M/f_1M$, $\overline{f}_2, \ldots, \overline{f}_r$
satisfy the conditions of the lemma, by induction we conclude
this complex is acyclic in postive degrees. This finishes the proof
of the first statement. The second statement immediately follows
from the first.
\end{proof}
